\documentclass{article}
\usepackage{mathtools} 
\usepackage{fontspec}
\usepackage[UTF8]{ctex}
\usepackage{amsthm}
\usepackage{mdframed}
\usepackage{xcolor}
\usepackage{amssymb}
\usepackage{amsmath}


% 定义新的带灰色背景的说明环境 zremark
\newmdtheoremenv[
  backgroundcolor=gray!10,
  % 边框与背景一致，边框线会消失
  linecolor=gray!10
]{zremark}{说明}

% 通用矩阵命令: \flexmatrix{矩阵名}{元素符号}{行数}{列数}
\newcommand{\flexmatrix}[4]{
  \[
  #1 = \begin{pmatrix}
    #2_{11}     & #2_{12}     & \cdots & #2_{1#4}   \\
    #2_{21}     & #2_{22}     & \cdots & #2_{2#4}   \\
    \vdots      & \vdots      & \ddots & \vdots     \\
    #2_{#31}    & #2_{#32}    & \cdots & #2_{#3#4}
  \end{pmatrix}
  \]
}

% 简化版命令（默认矩阵名为A，元素符号为a）: \quickmatrix{行数}{列数}
\newcommand{\quickmatrix}[2]{\flexmatrix{A}{a}{#1}{#2}}

\begin{document}
\title{注释 7.2}
\author{张志聪}
\maketitle

\begin{zremark}
  7.7' (1)的证明
\end{zremark}

\textbf{证明：}

只证明充分性。

由题设和定理7.4可知，最大聚点是存在的。

设最大聚点为$A$，由对任意$\alpha > \overline{A}$，$\{x_n\}$
中大于$\alpha$的项至多有限个可知，$A \leq \overline{A}$。

假设$A < \overline{A}$，设$\epsilon = \frac{\overline{A} - A}{2}$，
于是可得$A + \epsilon < \overline{A}$，
由定理7.7(i)可知，$A + \epsilon$的右侧只有有限个项，
这与题设对任意$\beta < \overline{A}$，$\{x_n\}$中
大于$\beta$的项有无限多个相矛盾。

\end{document}